\section{De la AI Bubble a la AI War: por qué la «burbuja» no basta como explicación}

Esta sección aborda primero el marco general de todo el episodio. Al inicio del programa la pregunta es por la «AI Bubble», pero la respuesta de Guang Mi (广密) es que «se parece más a una AI War». No es un juego de palabras, sino un cambio en la unidad de análisis. Si la IA se ve solo como un proyecto comercial, se parte del flujo de efectivo, la valuación y el gasto de capital, y se pregunta «cómo se recupera este dinero»; si la IA se ve como la siguiente generación de plataformas, como capacidad de defensa nacional y como una reorganización del orden tecnológico, también hay que preguntar «si no se invierte, ¿se corre el riesgo de quedar desplazado?».

\begin{figure}[H]
\centering
\includegraphics[width=0.96\textwidth]{ai-war-alliances.png}
\caption{Las dos líneas principales de la AI War: el ecosistema de GPU de Nvidia, el ecosistema de TPU de Google y las alianzas de contrapeso que se forman en torno a OpenAI/Google. Diagrama conceptual de elaboración propia, basado en la conversación de 00:15:04--00:17:48.}
\end{figure}

\begin{knowledgebox}{Lectura de la figura: primero los ecosistemas, luego las alianzas}
Arriba a la izquierda y arriba al centro están, respectivamente, las dos líneas principales de ecosistemas de hardware: las GPU de Nvidia y las TPU de Google. OpenAI, Anthropic y otras empresas líderes de modelos se ubican sobre todo en el ecosistema de GPU; Google, en cambio, integra de extremo a extremo modelos, chips, nube y productos. Abajo, la «alianza anti-Google» y la «alianza anti-OpenAI» recuerdan que cuanto más fuerte es un actor individual, más contrapesos provoca por parte de otros gigantes, del capital y de los socios del ecosistema.
\end{knowledgebox}

\subsection{Primer sentido de la AI War: defensa comercial}

Antes se dijo que la AI War no es una figura retórica; este apartado examina primero por qué, en el plano comercial, nadie puede permitirse perder. Guang Mi da dos ejemplos concretos. El primero es que los proveedores de nube quedan reducidos a marca blanca frente a los modelos: antes, desarrolladores y empresas consideraban primero AWS, Azure o Google Cloud; ahora muchas aplicaciones nuevas preguntan primero «qué modelo usar»: GPT, Claude o Gemini. La nube sigue siendo importante, pero la atención de los usuarios podría desplazarse de la nube al modelo.

El segundo es que el Agent podría convertirse en un nuevo punto de entrada del tráfico. La búsqueda tradicional, las Super App, las OTA, las plataformas de transporte, de empleo y de comercio electrónico construyen muchos de sus modelos de negocio sobre la distribución de información, la intermediación y las comisiones por transacción. Si un Agent puede comparar, negociar, hacer pedidos, pagar y dar seguimiento en nombre del usuario, las comisiones y el control del punto de entrada de las plataformas tendrán que valorarse de nuevo. Para los gigantes no se trata de «una herramienta más», sino de una batalla defensiva por fuentes de utilidades que ya tienen.

\begin{warningbox}{Advertencia de interpretación: AI War no significa que las cuentas de corto plazo cuadren}
Decir que la IA es una guerra no implica que cada proyecto de centro de datos, cada plan de entrenamiento de modelos y cada startup merezcan inversión. Solo indica que la conducta de los actores clave no puede explicarse con el ciclo de retorno de una inversión SaaS común, porque el costo de oportunidad de «no invertir» puede ser que se reescriba su posición de plataforma.
\end{warningbox}

\subsection{Segundo sentido de la AI War: estrategia nacional y movilización de capacidad de cómputo}

Guang Mi va más allá y compara la IA con una nueva forma de defensa nacional y con las armas nucleares. Esta frase puede oírse como una exageración, pero apunta a la estructura de capacidades a escala nacional: quien tenga modelos más fuertes, más capacidad de cómputo, una mayor concentración de investigadores y una cadena de suministro de chips más completa podría obtener una ventaja sistémica en la producción de software, el descubrimiento científico, el apoyo militar, la automatización industrial y la distribución de información.

Por eso proyectos cuyas cuentas de flujo de efectivo no cuadran a corto plazo siguen recibiendo apoyo del capital y de las políticas públicas. Las empresas temen quedar desplazadas, los Estados temen rezagarse estratégicamente, y los mercados de capital buscan las pocas direcciones capaces de absorber enormes sumas sin dejar de ofrecer expectativas de crecimiento. Con estos tres factores superpuestos, el ritmo del gasto de capital en IA difícilmente se explica solo con la palabra «burbuja».

\begin{importantbox}{Las tres cuentas para analizar la inversión en IA}
\begin{tabularx}{\textwidth}{p{0.20\textwidth}p{0.35\textwidth}X}
\toprule
Cuenta & Pregunta central & Juicio correspondiente en este episodio \\
\midrule
Cuenta financiera & Cómo cubren la inversión los ingresos, el margen bruto, el flujo de efectivo y la depreciación & El compromiso de 1.4T de OpenAI será cuestionado a corto plazo. \\
Cuenta de plataforma & Si el nuevo punto de entrada sustituirá al anterior y quién controla la relación con el usuario & El Agent y la nube de modelos podrían reescribir la búsqueda, la nube y las Super App. \\
Cuenta estratégica & Si, al no invertir, el país y la empresa pierden capacidades futuras & La IA se compara con la defensa nacional y las armas nucleares; los actores no se atreven a retirarse a la ligera. \\
\bottomrule
\end{tabularx}
\end{importantbox}

\subsection{Del debate sobre la burbuja al análisis del panorama competitivo}

El debate sobre la burbuja suele preguntar «si la valuación es demasiado alta». La pregunta que más vale conservar de este episodio es «quién se verá obligado a seguir invirtiendo, quién puede formar alianzas y quién logra convertir la inversión en un ciclo cerrado de producto y datos». Ese es también el orden de lo que sigue: primero se desglosa la cuenta de ingresos de OpenAI, luego se examinan los bandos de Nvidia/Google, después se discute por qué los tres grandes modelos se alternan el liderazgo y, por último, se llega al Online Learning, una variable que podría cambiar el panorama actual.

\subsection{Resumen de la sección}

El marco de la AI War devuelve la discusión de los ánimos sobre la valuación a las variables estructurales. No niega que haya burbujas parciales, pero recuerda al lector que, cuando la IA amenaza al mismo tiempo a la nube, la búsqueda, las Super App, el trabajo de oficina y la competitividad nacional, los gigantes y los Estados invertirán de forma más agresiva que en un proyecto comercial común.
